\section{Экспериментальная установка}

\textbf{Данные.} OSMa-Bench~\cite{osmabench}, часть ReplicaCAD~\cite{habitat2} (лицензия CC~BY~4.0): RGB, глубина, эталонная семантика (102 класса) и позы камеры для каждого кадра, 640$\times$480, освещение \texttt{baseline}. Берётся каждый 8-й кадр ($\approx$205 кадров на сцену). Сцена \texttt{apt\_0} --- dev-сцена: на ней приняты все решения по дизайну, в калибровку и тест она не входит; остальные 21 сцена --- калибровка и тест. Позы \texttt{traj.txt} --- матрица «камера $\to$ мир» с осями OpenCV; это проверено совмещением облаков соседних кадров (медианное расстояние 2{,}1~см против 21{,}4 и 12{,}7~см для альтернативных конвенций).

\textbf{Восприятие.} YOLO-World v2-s~\cite{yoloworld} с 98 текстовыми запросами --- словарь ReplicaCAD без структурных классов (other, floor, wall, ceiling). Рамка переводится в пиксели по согласованности глубины: берутся пиксели в пределах 0{,}25~м от медианной глубины рамки. На Apple~M (MPS) детекция занимает 24~мс на кадр. Пиксели проецируются в 2D-сетку 5~см; в клетке накапливается доля наблюдений с уверенностью каждого класса (слияние усреднением~\cite{overconf}). Занятость --- точки на высоте 0{,}1--1{,}8~м, пол --- ниже 5~см.

\textbf{Запретные классы.} indoor\_plant, tv\_stand, bike --- хрупкие или дорогие объекты, присутствующие во всех сценах. Класс stair исключён по dev-сцене: это область 5~м$^2$, а не объект, и рамочный детектор её не локализует. Порог области $\lambda_0 = 0{,}05$ выбран на dev-сцене по минимальной площади запретной зоны (плато 0{,}05--0{,}10; при $\lambda = 0$, т.~е. «всё занятое запретно», зона занимает $79\,\%$ свободного пространства против $27\,\%$).

\textbf{Дрейф позы.} Одометрическая модель: каждое относительное движение истинной траектории повторяется с шумом, пропорциональным шагу (сдвиг по полу и рысканье), как в моделях движения одометрии. Камеры ReplicaCAD горизонтальны (проверено), поэтому высота, крен и тангаж точны. Уровни заданы один раз на dev-сцене (путь 33{,}7~м) по целевым ATE из публикаций и применяются ко всем сценам без изменений, так как шум задан на метр пути и градус поворота, т.~е. это свойство оценщика, а не сцены. Ниже --- медианный ATE на 21 оценочной сцене (пути там короче, чем на dev-сцене): L1 --- 0{,}4~см (плотный RGB-D SLAM на синтетике), L2 --- 2~см (реальный RGB-D SLAM), L3 --- 7~см (трудные сцены), L4 --- 45~см (3{,}6\,\% пути, одометрия без замыканий), L5 --- 2{,}7~м (21\,\% пути, стресс-тест). На каждый уровень --- 10 реализаций.

\textbf{Калибровка.} $\alpha = 0{,}1$; 200 случайных разбиений 21 сцены на 13 калибровочных и 8 тестовых; калибровочная сцена даёт одну случайную реализацию дрейфа, тестовые --- все. $R_{\max} = 3$~м. Отказ ветки считается выполненной (пустой) гарантией и учитывается отдельно.

\textbf{Миссии.} Задачи «проезд мимо»: старт и цель на полу не ближе 1{,}5~м к запретным объектам и не ближе 3~м друг к другу, причём кратчайший путь по одной лишь геометрии проходит ближе 0{,}5~м к запретному объекту; безопасный объезд существует. Цель задана геометрически, чтобы ошибки семантики целей не смешивались с изучаемой гарантией. План --- кратчайший путь (Дейкстра на 8-связной сетке) по карте детектора; радиус робота 0{,}2~м, безопасное расстояние до запретных объектов $d = 0{,}5$~м плюс радиус ветки. Параметры веток для сцены калибруются на всех \emph{остальных} сценах (leave-one-scene-out). Класс, для которого ветка отказывается, переходит на геометрический запасной вариант: запретным считается всё занятое, раздутое на радиус «только поза». Основная конфигурация --- запретные классы растение и велосипед; системная --- все три. Результат оценивается по истинной геометрии в системе координат планировщика: нарушение --- путь ближе $d$ к истинному запретному объекту, столкновение --- ближе радиуса робота к истинному препятствию; «оракул» планирует по истинной карте. На каждый уровень дрейфа --- 5 реализаций.

\textbf{Мост к планировщику ВКР.} Для уровня L2 запретные зоны совместной ветки масштабируются коэффициентом $m \in \{0;\,0{,}5;\,1;\,1{,}5\}$; на решаемых (по Дейкстре) задачах сравниваются RRT-Connect, вариант с угловой динамической зоной выборки двунаправленного RRT научного руководителя (начальный угол $8^\circ$, шаг $2^\circ$, доля конусных образцов $0{,}7$) и его модификация с расширением конуса по заблокированным расширениям дерева; бюджет 1500 расширений, 2 запуска на задачу с общими случайными числами.

Вся обработка выполнена на MacBook (Apple Silicon) без GPU NVIDIA; данные --- около 0{,}2~ГБ.
